% ==========================================
% CHAPTER 6
% ==========================================
\chapter{Functional Programming} % 

\section{Variable Scope in Functions} % 
As you begin to write longer, more complex code, the ability to organize it into individual sub-components becomes increasingly useful. %  
The primary way we do this in Python is via functions and classes. % 

We can write a function to make the loading of FITS images into our code a few lines faster later on: % 

\begin{lstlisting}
def load_img(fname, extension=0):
    with fits.open(fname) as hdulist:
        img = hdulist[extension].data
    return img
\end{lstlisting} % 

An important aspect of functions is that the variables defined and used within a function are what is known as ``local in scope.'' %  
That means that those variables are created when the function is called and destroyed once the output is returned. % 

One important way to ensure we are doing all this correctly is good documentation. %  
The way to set documentation for a function is to place it inside triple quotes right at the top of our function definition: % 

\begin{lstlisting}
def load_img(fname, extension=0):
    """
    A function to quickly extract the data extension of a fits file
    INPUTS:
    fname (string): path/file name of the fits file to be loaded
    extension (int) [default: 0]: extension to index
    RETURNS:
    img (array_like): data attribute queried
    """
    with fits.open(fname) as hdulist:
        img = hdulist[extension].data
    return img
\end{lstlisting} % 

When you have set documentation this way, you can run the command \texttt{help(load\_img)} from the Python/iPython interpreter. % 

\section{Setting optional arguments, args, and kwargs} % 
\subsection{Optional Arguments} % 
It's possible that you may have a function that can take many potential arguments, but most uses of the function will have the same defaults for many arguments save a few critical ones. %  

\begin{lstlisting}
def somefunction(var1, var2, var3=1, var4='cat'):
    output = str(var1+var2+var3) + var4
    return output
\end{lstlisting} % 

\subsection{Args and Kwargs} % 
The \texttt{*args} and \texttt{**kwargs} commands allow us to feed variable numbers of arguments to a function. % 

\begin{lstlisting}
def test_function(farg, *args):
    print('formal argument:', farg)
    for arg in args:
        print('new arg:', arg)
\end{lstlisting} % 

If those are \texttt{*args}, what about \texttt{**kwargs}? %  
Keyword-args let you pass a variable number of extra variables to the function. %  
The difference is, when you feed those extra arguments into the function, you individually give each a new keyword by setting it equal to it in the function call. % 

\begin{lstlisting}
def kwarg_examplefunction(farg, **kwargs):
    print('formal argument:', farg)
    for key in kwargs:
        print('argument:', kwargs[key])
\end{lstlisting} % 

The way, in closing, to call this function would be: % 

\begin{lstlisting}
kwarg_examplefunction('tree', arg1='cat', blah='dog')
\end{lstlisting} % 